\chapter{Setup \& Directory Tour}
\label{ch:setup-directory-tour}

\textit{This chapter walks you through installing \sysname{}'s dependencies,
configuring the environment, and starting the infrastructure for the first time.
By the end you will have a running system and a mental map of where everything
lives in the codebase.}

% ─────────────────────────────────────────────────────────────────────
\section{Prerequisites}
\label{sec:prerequisites}

Before cloning the repository, ensure the following tools are installed on your
machine:

\begin{itemize}
    \item \textbf{Node.js 24 LTS} --- the frontend runtime.  Install via
          \ic{nvm install 24} or from \url{https://nodejs.org}.
    \item \textbf{Python 3.12+} --- the backend runtime.  Verify with
          \ic{python3 --version}.
    \item \textbf{Docker \& Docker Compose} --- runs PostgreSQL, Neo4j, and
          Redis in containers.  Docker Desktop includes Compose on macOS and
          Windows.
    \item \textbf{pnpm 9+} --- the monorepo package manager.  Install with
          \ic{npm install -g pnpm} or \ic{corepack enable}.
    \item \textbf{NVIDIA NIM API key} --- obtain a free key from
          \url{https://build.nvidia.com}.  You need \emph{two} keys: one for the
          LLM (\confkey{NVIDIA\_API\_KEY}) and one for the embedding model
          (\confkey{NVIDIA\_EMBEDDING\_API\_KEY}).  Alternatively, configure
          AWS credentials for Amazon Bedrock.
\end{itemize}

\begin{keyinsight}
\sysname{} uses two separate NVIDIA API keys because the LLM and embedding
endpoints are billed independently.  You can generate both from the same
\texttt{build.nvidia.com} account.
\end{keyinsight}


% ─────────────────────────────────────────────────────────────────────
\section{Environment Setup}
\label{sec:environment-setup}

Follow these six steps in order.  Each step is idempotent---running it twice
will not break anything.

\medskip
\noindent\textbf{Step 1.  Clone the repository.}

\begin{lstlisting}[style=bash, numbers=none]
git clone https://github.com/shehral/continuum.git && cd continuum
\end{lstlisting}

\noindent\textbf{Step 2.  Create the environment file.}

\begin{lstlisting}[style=bash, numbers=none]
cp .env.example .env
\end{lstlisting}

Open \codefile{.env} in your editor and set the following variables.  The
template ships with placeholder values that \emph{must} be replaced.

\begin{table}[ht]
\centering
\caption{Key environment variables to configure.}
\label{tab:env-vars}
\begin{tabularx}{\textwidth}{l X}
\toprule
\textbf{Variable} & \textbf{What to set} \\
\midrule
\confkey{NVIDIA\_API\_KEY}           & Your NVIDIA NIM LLM key (starts with \ic{nvapi-}). \\
\confkey{NVIDIA\_EMBEDDING\_API\_KEY} & Your NVIDIA NIM embedding key. \\
\confkey{POSTGRES\_PASSWORD}         & A strong random password.  Generate one with
                                       \ic{python3 -c "import secrets; print(secrets.token\_urlsafe(24))"}. \\
\confkey{NEO4J\_PASSWORD}            & Another random password for the graph database. \\
\confkey{REDIS\_PASSWORD}            & Another random password for the cache. \\
\confkey{DATABASE\_URL}              & Update the password segment to match
                                       \confkey{POSTGRES\_PASSWORD}. \\
\confkey{REDIS\_URL}                 & Update the password segment to match
                                       \confkey{REDIS\_PASSWORD}. \\
\confkey{SECRET\_KEY}                & A random 32-byte secret.  Generate with
                                       \ic{python3 -c "import secrets; print(secrets.token\_urlsafe(32))"}. \\
\confkey{NEXTAUTH\_SECRET}           & \textbf{Must be identical to} \confkey{SECRET\_KEY}. \\
\bottomrule
\end{tabularx}
\end{table}

\begin{warning}
\confkey{SECRET\_KEY} and \confkey{NEXTAUTH\_SECRET} must be the \emph{same}
string.  The backend signs JWTs with \confkey{SECRET\_KEY}; the frontend
verifies them with \confkey{NEXTAUTH\_SECRET}.  A mismatch causes silent
authentication failures---you will be able to register but every subsequent
request falls back to ``anonymous.''
\end{warning}

\noindent\textbf{Step 3.  Start infrastructure.}

\begin{lstlisting}[style=bash, numbers=none]
docker-compose up -d
\end{lstlisting}

This launches three containers, all bound to \ic{127.0.0.1} (localhost only):

\begin{itemize}
    \item \textbf{PostgreSQL 18} on port \ic{5432} --- relational storage for
          users, projects, and decision metadata.
    \item \textbf{Neo4j 2025.01} on ports \ic{7474} (HTTP browser) and
          \ic{7687} (Bolt protocol) --- the knowledge graph.
    \item \textbf{Redis 7.4} on port \ic{6379} --- caching, rate limiting, and
          session storage.
\end{itemize}

\noindent\textbf{Step 4.  Install frontend dependencies.}

\begin{lstlisting}[style=bash, numbers=none]
pnpm install
\end{lstlisting}

\noindent\textbf{Step 5.  Set up the backend.}

\begin{lstlisting}[style=bash, numbers=none]
cd apps/api && python3 -m venv .venv \
  && .venv/bin/pip install -e ".[dev]" && cd ../..
\end{lstlisting}

This creates an isolated Python virtual environment and installs all backend
dependencies including development tools (pytest, ruff).

\noindent\textbf{Step 6.  Run database migrations.}

\begin{lstlisting}[style=bash, numbers=none]
pnpm db:migrate
\end{lstlisting}

Under the hood this runs \ic{alembic upgrade head} inside the virtual
environment, creating all PostgreSQL tables and indices.


% ─────────────────────────────────────────────────────────────────────
\section{Directory Tour}
\label{sec:directory-tour}

The project follows a monorepo structure managed by \ic{pnpm-workspace.yaml}.
Below is the top-level layout with annotations.

\begin{figure}[ht]
\centering
\begin{forest}
  for tree={
    font=\small\ttfamily,
    grow'=0,
    child anchor=west,
    parent anchor=south,
    anchor=west,
    calign=first,
    edge path={
      \noexpand\path [draw, -{Stealth[length=3pt]}]
        (!u.south west) +(7.5pt,0) |- (.child anchor)\forestoption{edge label};
    },
    before typesetting nodes={
      if n=1 {insert before={[,phantom]}} {}
    },
    fit=band,
    before computing xy={l=15pt},
  }
[continuum/
  [apps/
    [web/ \textnormal{\scriptsize\color{ContinuumGray} --- Next.js 16 frontend}]
    [api/ \textnormal{\scriptsize\color{ContinuumGray} --- FastAPI backend}
      [routers/ \textnormal{\scriptsize\color{ContinuumGray} --- 13 API endpoint modules}]
      [services/ \textnormal{\scriptsize\color{ContinuumGray} --- 14 business logic modules}]
      [models/ \textnormal{\scriptsize\color{ContinuumGray} --- SQLAlchemy + Pydantic schemas}]
      [evaluation/ \textnormal{\scriptsize\color{ContinuumGray} --- Benchmark framework}]
      [middleware/ \textnormal{\scriptsize\color{ContinuumGray} --- Security, rate limiting, metrics}]
      [db/ \textnormal{\scriptsize\color{ContinuumGray} --- Database connections (PG, Neo4j, Redis)}]
      [tests/ \textnormal{\scriptsize\color{ContinuumGray} --- Test suite}]
    ]
    [mcp/ \textnormal{\scriptsize\color{ContinuumGray} --- MCP server tools}]
  ]
  [docker-compose.yml \textnormal{\scriptsize\color{ContinuumGray} --- PG 18, Neo4j 2025.01, Redis 7.4}]
  [infra/ \textnormal{\scriptsize\color{ContinuumGray} --- Grafana \& Prometheus dashboards}]
  [pnpm-workspace.yaml \textnormal{\scriptsize\color{ContinuumGray} --- Monorepo workspace config}]
]
\end{forest}
\caption{Top-level project structure.  The \codefile{apps/api/} directory
contains the bulk of the backend logic.}
\label{fig:directory-tree}
\end{figure}

The split is deliberate: the \codefile{apps/web/} directory is a self-contained
Next.js application that communicates with the backend exclusively through
HTTP and WebSocket endpoints.  The \codefile{apps/api/} directory is a
standalone FastAPI application with its own virtual environment, configuration,
and test suite.  The \codefile{apps/mcp/} directory provides a Model Context
Protocol server that allows AI agents to query and write to the knowledge graph.


% ─────────────────────────────────────────────────────────────────────
\section{Key Files to Bookmark}
\label{sec:key-files}

As you work through this guide, you will return to certain files repeatedly.
Table~\ref{tab:key-files} lists the most important ones and when to reach for
them.

\begin{table}[ht]
\centering
\caption{Files you will reference frequently.}
\label{tab:key-files}
\begin{tabularx}{\textwidth}{l l X}
\toprule
\textbf{File} & \textbf{Purpose} & \textbf{Read this when\ldots} \\
\midrule
\codefile{config.py}          & All settings           & \ldots you need to change any configuration. \\
\codefile{entity\_resolver.py} & Core resolution algorithm & \ldots you are working on entity resolution. \\
\codefile{graph\_rag.py}       & Retrieval pipeline     & \ldots you are working on search or the Ask feature. \\
\codefile{ontology.py}        & Canonical mappings     & \ldots you are adding or modifying entity mappings. \\
\codefile{extractor.py}       & Decision extraction    & \ldots you are modifying extraction prompts. \\
\codefile{llm.py}             & LLM client             & \ldots you are debugging LLM calls or switching providers. \\
\codefile{embeddings.py}      & Embedding client       & \ldots you are working on vector search. \\
\bottomrule
\end{tabularx}
\end{table}

All paths in Table~\ref{tab:key-files} are relative to \codefile{apps/api/services/},
except \codefile{config.py} (which lives at \codefile{apps/api/config.py}) and
\codefile{ontology.py} (which lives at \codefile{apps/api/models/ontology.py}).


% ─────────────────────────────────────────────────────────────────────
\section{Common First-Run Problems}
\label{sec:first-run-problems}

If things go wrong during setup, check these common causes first.

\begin{warning}
\textbf{Port conflicts.}  If another service is already listening on port
\ic{5432}, \ic{7474}, or \ic{6379}, Docker will fail to bind.  Stop the
conflicting service or edit \codefile{docker-compose.yml} to map to a different
host port.  On macOS, \ic{lsof -i :5432} will show what is using the port.
\end{warning}

\begin{warning}
\textbf{Alembic \texttt{DuplicateTableError}.}  This happens when the database
already has tables but Alembic's version tracking table is missing (for example,
after restoring a database dump).  The fix is to stamp the current state without
re-running migrations:

\begin{lstlisting}[style=bash, numbers=none]
cd apps/api && .venv/bin/alembic stamp head
\end{lstlisting}
\end{warning}

\begin{warning}
\textbf{Silent authentication failure.}  You can register and log in, but every
API request returns data as if you were anonymous.  This almost always means
\confkey{NEXTAUTH\_SECRET} does not match \confkey{SECRET\_KEY}.  After
correcting the mismatch, restart both the frontend and backend.
\end{warning}

\begin{tryit}
Run \ic{docker-compose up -d} and verify that Neo4j is accessible at
\url{http://localhost:7474}.  The default credentials are
\ic{neo4j} / the password you set in \confkey{NEO4J\_PASSWORD}.  You should see
the Neo4j Browser interface with an empty database ready for use.
\end{tryit}
